\section{Sections, scales, and a quantitative separation}
\label{sc:section}

Throughout this section the manifold and the density bounds are fixed.
We use the negative-cost convention
\[
 f^c(x)=\max_{y\in M}\{-c(x,y)-f(y)\}.
\]
Thus a cost-dual pair satisfies $u=v^c$, $v=u^c$, and
$u(x)+c(x,y)+v(y)\geq0$.
Let $\mathcal C$ denote the class of these pairs for all admissible density
pairs, with any fixed additive normalization.
All members of $\mathcal C$ have a common Lipschitz bound, and their
oscillations are uniformly bounded.
None of the arguments below uses derivatives of the densities.

\subsection{A scalar quantity controlling both transport directions}

For $(u,v)\in\mathcal C$ and $r\geq0$, put
\[
 S_r(x;u,v)=\{y\in M:u(x)+c(x,y)+v(y)\leq r\},
 \qquad
 F(r)=\sup_{(u,v)\in\mathcal C,\ x\in M}
       \osc_{S_r(x;u,v)}v.
\]
The sets $S_r$ are nonempty and compact.  The function $F$ is bounded
and nondecreasing.

\begin{lemma}[The power reduction]\label{sc:power}
If there are $\beta>0$, $C_0<\infty$, and $r_0>0$ such that
$F(r)\leq C_0r^\beta$ for $0<r<r_0$, then all the optimal maps in the
fixed density class have homeomorphic representatives with a common
H\"older estimate in both directions.
One may take any exponent no larger than
$\frac12\min\{\beta,1\}$.
\end{lemma}

\begin{proof}
Write
\[
 h_x(p)=u(x)+\frac12|p|^2+v(\exp_xp),\qquad p\in G_x.
\]
By Corollary~\ref{geo:sections}, the lifted section
$\{p\in G_x:h_x(p)\leq r\}$ is convex.
If $p,q$ belong to it and $m=(p+q)/2$, the squared-norm identity gives
\begin{align}
 \frac18|p-q|^2
 &=v(\exp_xm)-\frac{v(\exp_xp)+v(\exp_xq)}2
   +\frac{h_x(p)+h_x(q)}2-h_x(m) \notag\\
 &\leq F(r)+r.                                      \label{sc:diameter}
\end{align}
The exponential map has a common Lipschitz bound on the compact bundle
of minimizing velocities.  Hence the target diameter of a section is
at most $C\sqrt{F(r)+r}$.

Let $(x,y)$ and $(x',y')$ be contacts of the same dual pair.  The common
Lipschitz bounds for $u$ and $c$ imply
\[
 0\leq u(x)+c(x,y')+v(y')\leq C d(x,x').
\]
Thus $y,y'$ belong to a section at $x$ of that height.  For sufficiently
small $d(x,x')$, \eqref{sc:diameter} gives
\[
 d(y,y')\leq C d(x,x')^{\frac12\min\{\beta,1\}}.
\]
In particular, letting the section height tend to zero shows that
every pole has exactly one contact endpoint.  The bound proves that
this endpoint depends continuously on the pole. The class also contains
the reversed pair $(v,u)$ after renormalization, which does not affect
its section oscillations. The same reasoning therefore gives a unique
continuous contact in the reverse direction. The dual identities give a contact above every point
of either copy of $M$, and the two resulting maps are inverse.
They agree almost everywhere with the optimal maps by optimality and
uniqueness.  Finally, compactness absorbs the remaining, larger
distances into the same uniform constant.
\end{proof}

We shall contradict failure of a power bound for $F$.
The following elementary selection makes the required quantifiers
explicit.

\begin{lemma}[Logarithmic plateaus]\label{sc:plateau}
Suppose that $F$ has no power upper bound near zero.
There are $r_j<R_j\to0$, pairs $(u_j,v_j)\in\mathcal C$, poles $x_j$,
and numbers $a_j,D_j$ with $D_j>0$ such that
\begin{gather}
 \min_{S_{r_j}(x_j;u_j,v_j)}v_j=a_j,\qquad
 \max_{S_{r_j}(x_j;u_j,v_j)}v_j=a_j+D_j,               \label{sc:endlevels}\\
 \frac{r_j}{R_j}\longrightarrow0,\qquad
 F(R_j)\leq(1+\delta_j)D_j,\qquad \delta_j\longrightarrow0,\qquad
 \frac{|\log D_j|}{|\log R_j|}\longrightarrow0.       \label{sc:plateau-data}
\end{gather}
Consequently $R_j/D_j^k\to0$ for every fixed $k>0$.
\end{lemma}

\begin{proof}
If $F$ vanished at a positive scale, monotonicity would already give a
power bound below that scale.  Thus $F(r)>0$.
Fix $C_F\geq\max\{1,\sup_r F(r)\}$ and set
\[
 f(s)=\log\frac{C_F}{F(e^{-s})}.
\]
This function is nonnegative and nondecreasing.  Failure of every
power bound implies $\liminf_{s\to\infty}f(s)/s=0$; otherwise
$f(s)\geq\beta s$ eventually for some $\beta>0$.
Choose $N_j\to\infty$ with $f(N_j)=o(N_j)$, and divide
$[N_j/2,N_j]$ into
\[
 m_j=\left\lceil\sqrt{N_j\max\{1,f(N_j)\}}\right\rceil
\]
equal intervals.  Their length $N_j/(2m_j)$ tends to infinity.
Since their total $f$-increment is at most $f(N_j)$, some interval
$[A_j,B_j]$ has increment at most $f(N_j)/m_j=o(1)$.
Put $R_j=e^{-A_j}$ and $r_j=e^{-B_j}$.  Then
\[
 \log(R_j/r_j)=B_j-A_j\longrightarrow\infty,\qquad
 \frac{F(R_j)}{F(r_j)}=e^{f(B_j)-f(A_j)}\longrightarrow1.
\]
Choose an actual section with oscillation
$D_j\geq(1-1/j)F(r_j)$, discarding the first few indices.
Its extrema are attained, which gives \eqref{sc:endlevels}.
Moreover $A_j\geq N_j/2$ and $f(B_j)\leq f(N_j)=o(N_j)$,
so $|\log D_j|=o(A_j)$.
This proves \eqref{sc:plateau-data}.  Finally
\[
 \log(R_j/D_j^k)=-A_j-k\log D_j=-A_j+o(A_j)\longrightarrow-\infty.
\]
\end{proof}

We suppress the index $j$ from now on.  All asymptotic statements refer
to the sequence in Lemma~\ref{sc:plateau}, after the templates below
have been fixed.  It is not necessary that $D$ tend to zero.
Set
\begin{equation}
 b_-=r(R/r)^{1/3},\qquad b_+=r(R/r)^{2/3}.
 \label{sc:binterval}
\end{equation}
Then $r\ll b_-\ll b_+\ll R$.  Uniformly for $b\in[b_-,b_+]$ and
$0\leq\tau\leq8b/D$,
\begin{equation}
 \frac rb\to0,\quad \frac bR\to0,\quad
 \frac b{D^k}\to0\ (k>0\text{ fixed}),\quad
 \frac\tau D\to0,\quad \frac{\tau^2}{b}\to0.
 \label{sc:smallness}
\end{equation}
For example, $\tau^2/b\leq64b/D^2$.
These estimates, rather than an additional restriction on $R/r$,
will control the Taylor errors.

\subsection{Templates and the order of choices}

At the normalized level $s=(v-a)/D$, we will use the two
modifications $v+bB_c(s)$ and $v+bB_o(s)$.
The first is the center datum; we compare it with the double
cost transform of the second. A positive excess in this
comparison will produce an interior maximum. We first fix
the center profile and a cap on the outer values, then choose
the outer curvature parameter $K$ and its profile, and only
afterward take the gap scales small. The construction below
makes that order explicit.

We use the fixed numerical choices
\begin{equation}
 \eps_0=2^{-10},\qquad c_0=2^{-10},\qquad c_1=c_0/4.
 \label{sc:numerical}
\end{equation}
The constants are deliberately smaller than needed.

The value constraints below serve the envelope comparison: the center
profile stays near one on the lower half of the selected oscillation,
whereas the outer profile vanishes at level one and places a barrier
on intermediate levels. We also prescribe strong outer concavity
where the outer value is bounded by a fixed cap. That derivative
condition will enter the endpoint lower tests in
Section~\ref{sec:jets}; the height bound below keeps the constants in
the switch and excess bounds independent of $K$.

\begin{lemma}[Smooth templates]\label{sc:templates}
There is a fixed smooth function $B_c:\R\to(-1,1)$ with bounded
derivatives such that
\[
 B_c'<0\text{ on }\R,\qquad
 B_c(s)\geq1-\eps_0\quad(s\leq1/2),\qquad
 B_c(s)<0\quad(s\geq3/4).
\]
Given a cap $M_0\geq2$ and $K\geq1$, there are
$0<\eta\leq1/128$ and a smooth bounded function $B_o\geq0$, with
bounded derivatives, satisfying
\begin{align}
 &B_o(s)\geq1\quad(s\leq1-\eta),\qquad
 B_o(1)=0,\qquad 1\leq B_o(0)\leq1+\eps_0,       \label{sc:outer-values}\\
 &B_o(s)\geq B_o(0)+1
       \quad(1/16\leq s\leq1-\eta),            \label{sc:barrier}\\
 &B_o'(s)>0,\quad B_o''(s)\leq-K
       \quad\text{if }-2\eta\leq s\leq1-\eta
                      \text{ and }B_o(s)\leq M_0.     \label{sc:curvature}
\end{align}
One can arrange $\|B_o\|_\infty\leq2M_0+8$, independently of $K$,
and make $\eta$ arbitrarily smaller if desired.
\end{lemma}

\begin{proof}
Choose $\gamma$ so large that $2/(1+e^{\gamma/8})\leq\eps_0$, and set
\[
 B_c(s)=1-\frac{2}{1+e^{-\gamma(s-5/8)}}.
\]
This has all the asserted properties.  In particular, $|B_c'|$ has a
positive minimum on every fixed compact interval.

For the outer function let
\[
 \beta(z)=\begin{cases}e^{-1/z},&z>0,\\0,&z\leq0,\end{cases}
 \qquad
 \chi(z)=\frac{\beta(z)}{\beta(z)+\beta(1-z)}.
\]
Thus $\chi$ is a smooth step, equal to zero on $(-\infty,0]$ and to
one on $[1,\infty)$.  Put
\[
 \kappa=K+1,\qquad A=M_0+3,\qquad
 L=128A\kappa,\qquad h_0=\frac1{128\kappa},\qquad
 0<\eta\leq\min\left\{\frac1{128},\frac{\eps_0}{24L}\right\}.
\]
Define
\begin{align*}
 w(s)&=(L-\kappa s)
       \chi\!\left(\frac{s+3\eta}{\eta}\right)
       \left[1-\chi\!\left(\frac{s-h_0}{h_0}\right)\right],\\
 J(s)&=1+\frac{\eps_0}2+\int_0^s w(q)\,dq,\\
 B_o(s)&=J(s)
       \left[1-\chi\!\left(\frac{s-(1-\eta)}{\eta}\right)\right].
\end{align*}
The function $w$ is nonnegative and supported in $[-3\eta,2h_0]$.
Also
\[
 \int_{-3\eta}^0w\leq3\eta(L+3\kappa\eta)
       \leq6\eta L\leq\eps_0/4,\qquad
 \int_0^{2h_0}w\leq2Lh_0=2A.
\]
It follows that
\[
 1+\eps_0/4\leq J(s)\leq1+\eps_0/2+2A.
\]
On $[-2\eta,h_0]$, one has $J'=L-\kappa s>0$ and
$J''=-\kappa$.  Moreover
\[
 J(h_0)-J(0)=A-\frac1{32768\kappa}>M_0+1.
\]
Since $J$ is nondecreasing, the condition $B_o(s)\leq M_0$ on
$[-2\eta,1-\eta]$ forces $s<h_0$; here $B_o=J$, proving
\eqref{sc:curvature}.  The function $J$ is constant beyond
$2h_0\leq1/64$, so \eqref{sc:barrier} also follows.
The final factor is identically one up to $1-\eta$, and decreases to
zero only in $(1-\eta,1)$.  This proves the remaining assertions,
including smoothness at all joins and the uniform height bound.
\end{proof}

The cap $M_0$ will be fixed independently of $K$ by
Lemma~\ref{sc:outward}; the choice
$M_0=24(n+1)/c_1$ suffices.  The center template is already fixed.
After fixing $M_0$, choose $K$ and the outer template, and only then
take the scales small.  Besides \eqref{sc:smallness}, this order gives
\begin{equation}
 \delta_j=o(\eta),\qquad
 \frac bD\max\{\|B_o'\|_\infty,\|B_c'\|_\infty\}\to0,\qquad
 \frac b{D^2}\max\{\|B_o''\|_\infty,\|B_c''\|_\infty\}\to0.
 \label{sc:template-smallness}
\end{equation}
In particular, both scalar modifications used below are increasing.
Constants in intermediate local estimates may depend on a fixed
template, but the constants $c_0,c_1$ and the cap do not depend on $K$.

\subsection{The first switch}

Fix a selected section with extreme endpoints $y_0,y_1$, so that
$v(y_0)=a$ and $v(y_1)=a+D$.  Choose minimizing logs
$p_0,p_1\in G_{x_0}$ to these endpoints, and put $P=|p_0|^2$.
For every target write
\[
 s_y=\frac{v(y)-a}{D},\qquad
 w_{o,b}(y)=v(y)+bB_o(s_y),\qquad
 w_{c,b}(y)=v(y)+bB_c(s_y),\qquad
 g_b=w_{o,b}^c.
\]
Let $H_o=\|B_o\|_\infty$.  Since $B_o\geq0$,
\begin{equation}
 u-bH_o\leq g_b\leq u.
 \label{sc:g-bounds}
\end{equation}

We will move the pole toward $y_0$. Initially the outer modification
lowers the level-zero mountain by about $b$, while leaving $y_1$
unpenalized. As long as mountains with $s_y\geq1-\eta$ dominate, their
minimizing logs are shorter than the log to $y_0$, so they grow more
slowly along this motion. The barrier excludes intermediate-level
winners, forcing a switch with both a low and a high active endpoint.
Their two logs will then provide a supported intermediate log with a
prescribed squared norm; the next subsection converts this construction
into a positive excess.

We record the uniform distance estimate used in the proof.
Splitting a minimizing geodesic at its midpoint gives a smooth upper
support for the squared-distance cost with a uniformly bounded
Hessian.  Consequently, for any minimizing log $q$ from $x$ to $y$,
\begin{equation}
 c(\exp_x z,y)\leq c(x,y)-\langle q,z\rangle+C_g|z|^2
 \quad\text{for }|z|\text{ sufficiently small},
 \label{sc:cost-upper}
\end{equation}
with one geometric constant and radius.  Along a smooth curve, the
one-sided derivative of $-c(\cdot,y)$ is the maximum of the pairings
with its minimizing logs.  In particular, whenever this restriction
is differentiable,
\begin{equation}
 \frac d{d\tau}[-c(x_\tau,y)]
       \leq d(x_\tau,y)|\dot x_\tau|.
 \label{sc:advance}
\end{equation}
These assertions hold also at a cut endpoint.  The same upper
derivative bound applies to an active mountain of a maximum: at a
time when the maximum is differentiable, a semiconvex active mountain
has the same derivative.  Indeed its one-sided derivatives straddle
that derivative in the opposite order to their semiconvex ordering,
and hence coincide.  The restrictions in question are Lipschitz, so
these differential inequalities can be integrated.

\begin{lemma}[First switch with absolute level control]\label{sc:switch}
For all sufficiently small selected scales, uniformly for
$b\in[b_-,b_+]$, there is
\begin{equation}
 \frac b{32D}\leq\tau_b\leq\frac{8b}D
 \label{sc:switch-time}
\end{equation}
such that, at $x_b=\exp_{x_0}(\tau_b p_0)$, the function $g_b$ has
active endpoints $y_-$ and $y_+$ with
\[
 -2\eta\leq s_{y_-}<1/16,\qquad s_{y_+}\geq1-\eta.
\]
The original endpoints $y_0,y_1$ and every active endpoint at this
pole lie in one original section $S_{C b}(x_b;u,v)$, where $C$ depends
only on the geometry and $H_o$.  In particular every active level is
in $[-\delta_j,1+\delta_j]$.
For every minimizing log $q_i$ to an active endpoint,
\begin{equation}
 |q_i|^2=P-2D s_{y_i}+O_{H_o}(\tau_b+b)
        =P-2D s_{y_i}+o(D).
 \label{sc:switch-norms}
\end{equation}
The constants in \eqref{sc:switch-time} are independent of $K$ and
of the height and derivatives of the outer template.
\end{lemma}

\begin{proof}
Set $x_\tau=\exp_{x_0}(\tau p_0)$ for
$0\leq\tau\leq T:=8b/D$, and let $\Pi_\tau$ be parallel transport
along this geodesic.  For small scales $T<1/2$.  The left cost is exact:
\[
 f_y(\tau):=-c(x_\tau,y)-v(y),\qquad
 f_0(\tau)=-\frac12(1-\tau)^2P-a,\qquad
 q_0(\tau)=(1-\tau)\Pi_\tau p_0.
\]
The initial gaps at $y_0,y_1$ lie in $[0,r]$, whence
\begin{equation}
 P-|p_1|^2=2D+O(r),\qquad
 |p_0-p_1|^2\leq8(D+r),\qquad
 P-\langle p_0,p_1\rangle\leq6D
 \label{sc:initial-logs}
\end{equation}
for small scales; the second inequality is \eqref{sc:diameter}
applied using the actual oscillation $D$ of this section.
In particular $P\geq D$.

First exclude the middle levels.  If $1/16\leq s_y\leq1-\eta$, put
\[
 h_y(\tau)=f_y(\tau)-bB_o(s_y)
                -f_0(\tau)+bB_o(0).
\]
Initially $h_y(0)\leq r-b$.  Wherever $h_y\geq-b/2$, the identity
relating values and squared distances gives, for every minimizing
log $q_y$ at the current pole,
\[
 |q_y|^2-|q_0|^2
   =-2Ds_y-2b(B_o(s_y)-B_o(0))-2h_y
   \leq-D/8-b.
\]
Equations~\eqref{sc:advance} and $\dot x_\tau=\Pi_\tau p_0$ imply
$h_y'\leq-D/16$ at almost every such time.
The positive part of $h_y+b/2$ is therefore nonincreasing and starts
at zero.  Thus $h_y\leq-b/2$ throughout $[0,T]$.
No middle-level mountain can be active.

Let $G_L$ and $G_R$ be the maxima of $f_y-bB_o(s_y)$ over the compact
sets $s_y\leq1/16$ and $s_y\geq1-\eta$, respectively.
Both are continuous, and the preceding paragraph shows
$g_b(x_\tau)=\max\{G_L(\tau),G_R(\tau)\}$.
Initially $G_L(0)\leq u(x_0)-b$ and
$G_R(0)\geq f_1(0)\geq u(x_0)-r$.
While $G_R>G_L$, every active endpoint is right and comparison with
the left modified mountain gives
\[
 |q_y|^2\leq|q_0|^2-2D(1-\eta)+2bB_o(0)
             \leq|q_0|^2-D.
\]
By rationalizing the difference of the square roots in
\eqref{sc:advance},
\[
 \frac d{d\tau}\{g_b(x_\tau)-f_0(\tau)\}\leq-D/2
 \quad\text{almost everywhere while }G_R>G_L.
\]
This difference starts at most $r$ and is always at least
$-bB_o(0)\geq-2b$.  It cannot obey the displayed inequality throughout
$[0,T]$.  There is therefore a first equality time $\tau_b\leq T$.
At that time both maxima are attained, and the middle-level exclusion
gives $s_{y_-}<1/16$.
Throughout the preceding interval, including its endpoint,
\begin{equation}
 -bB_o(0)\leq g_b(x_\tau)-f_0(\tau)\leq r.
 \label{sc:left-anchor}
\end{equation}

We next anchor the absolute levels.  By \eqref{sc:g-bounds} and
\eqref{sc:left-anchor},
$u(x_\tau)-f_0(\tau)\leq bH_o+r$.
The initial-cost upper support and \eqref{sc:initial-logs} give
\begin{equation}
 f_1(\tau)\geq
 f_0(0)-r+(P-6D)\tau-C_g\tau^2=:\mathcal L(\tau).
 \label{sc:benchmark}
\end{equation}
Since $f_0(\tau)=f_0(0)+P\tau-P\tau^2/2$, it follows that
$f_0-f_1\leq r+6D\tau+C\tau^2=O(b)$ on this interval.
Thus both $y_0,y_1$ lie in an original section of height $Cb$.
For any active endpoint $y$,
\[
 u(x_\tau)+c(x_\tau,y)+v(y)
      =u(x_\tau)-g_b(x_\tau)-bB_o(s_y)\leq bH_o.
\]
All these endpoints belong to the same section of height $Cb<R$.
Its oscillation is at most $(1+\delta_j)D$ and it contains levels
$a,a+D$.  Hence every active level lies in
$[-\delta_j,1+\delta_j]$, proving in particular the required lower
bound for $s_{y_-}$.

To prove the lower time estimate, keep the winning left endpoint
$y_-$ fixed, and put
\[
 k(\tau)=f_{y_-}(\tau)-bB_o(s_{y_-})-\mathcal L(\tau).
\]
One has $k(0)\leq2r-b<-b/2$, while
$k(\tau_b)\geq0$, since $g_b(x_{\tau_b})\geq f_1(\tau_b)$.
Let $\theta$ be the last time before $\tau_b$ at which $k=-b/2$.
Then $k\geq-b/2$ on $[\theta,\tau_b]$.
On this interval, using $s_{y_-}\geq-2\eta$ and $B_o\geq0$, the
inequality defining this last strip gives
\begin{align*}
 |q_{y_-}|^2-|q_0|^2
 &\leq4\eta D+2r+12D\tau+2C_g\tau^2+b\\
 &\leq D
\end{align*}
for all sufficiently small scales.  This uses only
\eqref{sc:smallness}; in particular $\tau/D\to0$.
Rationalizing once more gives
$f_{y_-}'\leq(1-\tau)P+2D$ almost everywhere in the strip.
Since $\mathcal L'=P-6D-2C_g\tau$, we have $k'\leq10D$ there
for small scales.  Therefore
\[
 b/2\leq k(\tau_b)-k(\theta)\leq10D(\tau_b-\theta),
\]
which implies the stated weaker lower bound $\tau_b\geq b/(32D)$.
No upper bound for the template or its derivatives occurs in this
last estimate.

Finally, if $y_i$ is active at the switch, the exact equality of its
modified value with $g_b$ gives
\[
 |q_i|^2=(1-\tau_b)^2P-2Ds_{y_i}
         -2bB_o(s_{y_i})-2\{g_b(x_b)-f_0(\tau_b)\}.
\]
Equation~\eqref{sc:left-anchor} and \eqref{sc:smallness} prove
\eqref{sc:switch-norms}.
\end{proof}

\subsection{Positive excess and an interior maximum}

Write $\widehat w_b=g_b^c=w_{o,b}^{cc}$ and
\[
 E_b(y)=w_{c,b}(y)-\widehat w_b(y),\qquad
 E(b)=\max_{y\in M}E_b(y).
\]
Order reversal of one cost transform and the double-transform
inequality give
\begin{equation}
 v\leq\widehat w_b\leq w_{o,b},\qquad
 E_b(y)\leq bB_c(s_y)\leq b.
 \label{sc:excess-upper}
\end{equation}

\begin{lemma}[Uniform positive excess]\label{sc:excess}
For all sufficiently small selected scales and all
$b\in[b_-,b_+]$,
\begin{equation}
 c_0b\leq E(b)\leq b,
 \label{sc:excess-bounds}
\end{equation}
with the constant in \eqref{sc:numerical}, independently of $K$.
\end{lemma}

\begin{proof}
Use the two active logs at the switch supplied by
Lemma~\ref{sc:switch}.  Equation~\eqref{sc:switch-norms} and their
level bounds show that their squared norms lie on opposite sides of
$P-D/2$.  By continuity along their segment, choose
$p\in\partial g_b(x_b)$ with
\[
 |p|^2=P-D/2.
\]
Theorem~\ref{geo:support} says that $p$ is minimizing and that
$y=\exp_{x_b}p$ is a contact endpoint for $g_b$.  Thus
\[
 E_b(y)=v(y)+bB_c(s_y)+c(x_b,y)+g_b(x_b).
\]
Minorizing $g_b$ by the left modified mountain gives the direct bound
\begin{equation}
 E_b(y)\geq D(s_y-1/4)+P\tau_b-P\tau_b^2/2
                  +b\{B_c(s_y)-B_o(0)\}.
 \label{sc:direct-excess}
\end{equation}
If $s_y>1/2$, the right side is at least $D/4-3b$, which is larger
than $c_0b$ for small scales.

It remains to consider $s_y\leq1/2$.  The original support inequality
at $x_0$ yields
\[
 v(y)-a\geq c(x_0,y_0)-c(x_0,y)-r.
\]
Apply \eqref{sc:cost-upper} at $x_b$ with backward displacement
$-\tau_b\Pi_{\tau_b}p_0$ to get
\[
 c(x_b,y)-c(x_0,y)
       \geq-\tau_b\langle p,\Pi_{\tau_b}p_0\rangle-C\tau_b^2.
\]
The corresponding cost difference for $y_0$ is exactly
$-P\tau_b+P\tau_b^2/2$.  Therefore
\begin{equation}
 \frac{E_b(y)}b\geq B_c(s_y)-B_o(0)
   +\bigl(P-\langle p,\Pi_{\tau_b}p_0\rangle\bigr)\frac{\tau_b}b
   -\frac rb-C\frac{\tau_b^2}b.
 \label{sc:geometric-excess}
\end{equation}
The norm choice gives
\[
 P-\langle p,\Pi_{\tau_b}p_0\rangle
   =\frac{P-|p|^2+|p-\Pi_{\tau_b}p_0|^2}{2}\geq D/4.
\]
The middle term in \eqref{sc:geometric-excess} is consequently at
least $1/128$.  The template difference is at least $-2\eps_0$,
and the last two terms tend to zero uniformly by
\eqref{sc:smallness}.  The numerical choices
\eqref{sc:numerical} give the asserted lower bound.
The upper bound is \eqref{sc:excess-upper}.
\end{proof}

To obtain an interior maximum, we now subtract a penalty from the
excess. A small linear term preserves its positive lower bound,
and a quadratic term makes the upper amplitude boundary unfavorable.
Define
\begin{equation}
 \pi(b)=\frac{c_0}4b+\frac{2b^2}{b_+}.
 \label{sc:penalty}
\end{equation}
In particular, $\pi'(b)\geq c_1$.

\begin{lemma}[Regularized interior maxima]\label{sc:maximum}
Fix a sufficiently small selected scale.  For every $t<1$ sufficiently
close to one, the continuous function
\[
 \Phi_t(b,z)=Q_tg_b(z)+Q_{1-t}w_{c,b}(z)-\pi(b),\qquad
 (b,z)\in[b_-,b_+]\times M,
\]
has a positive global maximum with $b$ in the interior of its interval.
Both spatial terms are $C^{1,1}$, and both minimizing points are unique.
If $x$ and $y$ are these minimizing points at a maximum, there is
$p\in\partial g_b(x)\cap G_x$ such that
\begin{equation}
 z=\exp_x(tp),\qquad y=\exp_xp,\qquad
 Q_tg_b(z)+Q_{1-t}w_{c,b}(z)=E_b(y)>0.
 \label{sc:aligned}
\end{equation}
Moreover $(x,y)$ is before cut.
\end{lemma}

\begin{proof}
At $b=b_+$, \eqref{sc:excess-bounds} makes $E(b)-\pi(b)$ negative.
At $b=b_-$, it is at most $b_-$.  At
$b_*=c_0b_+/8$, which belongs to the interval for small scales,
\[
 E(b_*)-\pi(b_*)\geq c_0^2b_+/16.
\]
Since $b_-/b_+\to0$, these inequalities give a strictly positive
interior maximum with a positive separation from both boundary values.

For a fixed scale and $b$ in its compact interval,
$Q_tg_b\to Q_1g_b=-\widehat w_b$ uniformly as $t\uparrow1$.
This follows directly from the uniform bound for
$|c/t-c|$.  The functions $w_{c,b}$ have a common Lipschitz constant
$L_*$ on that interval, and
\[
 w_{c,b}-\tfrac12L_*^2(1-t)\leq Q_{1-t}w_{c,b}\leq w_{c,b}.
\]
Thus $\Phi_t\to E_b-\pi$ uniformly in $(b,z)$, preserving the
positive interior maximum when $t$ is sufficiently close to one.

Proposition~\ref{geo:intermediate} gives the asserted regularity and unique
pole for $Q_tg_b$.  To check the other term, the scalar function
$q\mapsto q+bB_c((q-a)/D)$ has positive derivative, and its first two
derivatives are bounded at the fixed scales.  Composing a Lipschitz
semiconvex function with this function preserves semiconvexity: the
lower Taylor support for $v$, its Lipschitz bound, and the ordinary
second-order Taylor inequality for the scalar function give a common
quadratic lower support for $w_{c,b}$.
For completeness, a minimizer in $Q_\epsilon w_{c,b}(z)$ has
$d(z,y)\leq2L_*\epsilon$, by comparison with $y=z$.
For small $\epsilon$, all minimizers therefore lie in one convex
normal ball, where the Hessian of $c(z,\cdot)/\epsilon$ dominates
the negative semiconvexity bound.  Strict convexity gives uniqueness.
In one fixed coordinate chart, let $K_*$ be the semiconvexity constant
and let $\mu/\epsilon$ be a lower Hessian bound for the smooth cost
term.  The subgradient inequalities at the minimizers $y,y'$ for
nearby $z,z'$ and the mixed cost derivative bound give
\[
  (\mu/\epsilon-K_*)|y-y'|^2
       \leq (C/\epsilon)|z-z'|\,|y-y'|.
\]
For small $\epsilon$ this proves local Lipschitz dependence on $z$.
The envelope gradient formula then gives $C^{1,1}$ regularity.
This is uniform in $b$ on the fixed interval.

At a maximum, the two spatial gradients sum to zero.  Their envelope
formulas show that the geodesic from $x$ to $z$, continued for the
remaining time $1-t$ with the same velocity, reaches $y$.
The pole minimum gives $p\in\partial g_b(x)$.
By Theorem~\ref{geo:support}, this full vector is minimizing and its
endpoint is a support for $g_b$.  Splitting its energy at time $t$
proves \eqref{sc:aligned} exactly.

Finally put $l=1+bB_c'(s_y)/D\in(0,1)$.
The short final minimization gives its backward velocity in
$\partial w_{c,b}(y)=l\partial v(y)$; the equality follows from the
increasing scalar chain rule, which follows in both directions by
Taylor expansion and inversion of that scalar function.
Dividing by $l$ gives a minimizing velocity by
Theorem~\ref{geo:support}.  The backward velocity to $x$ is a strict
radial shortening of this vector, and is therefore before cut.
\end{proof}

\subsection{The parameter inequality for every representation}

\begin{lemma}[The parameter derivative]\label{sc:parameter}
At any maximum in Lemma~\ref{sc:maximum}, let
\[
 p=\sum_{i=1}^N m_i p_i,\qquad m_i>0,\qquad\sum_i m_i=1,
\]
be any finite representation by minimizing logs active for the displayed
representation $g_b=w_{o,b}^c$ at its minimizing pole $x$.
Thus $g_b(x)+c(x,\exp_xp_i)+w_{o,b}(\exp_xp_i)=0$.
Write $y_i=\exp_xp_i$ and
$s_i=(v(y_i)-a)/D$, and let $s=(v(y)-a)/D$ be the center level.
Then
\begin{equation}
 B_c(s)-\sum_i m_iB_o(s_i)\geq \pi'(b)\geq c_1.
 \label{sc:parameter-inequality}
\end{equation}
\end{lemma}

\begin{proof}
Fix the prefix time and the spatial maximum point $z$, and abbreviate
$A(b)=Q_tg_b(z)$ and $C(b)=Q_{1-t}w_{c,b}(z)$.
Decrease $b$ by $h>0$, and let $x_h$ be the new pole for $A(b-h)$.
Because $0\leq g_{b-h}-g_b\leq hH_o$, compactness and uniqueness of
the old pole imply $x_h\to x$.
The strict local quadratic growth in
Proposition~\ref{geo:intermediate} gives, in normal coordinates
$\Delta_h=\exp_x^{-1}x_h$,
\begin{equation}
 |\Delta_h|=O(\sqrt h),\qquad 0\leq A(b-h)-A(b)\leq hH_o.
 \label{sc:pole-movement}
\end{equation}

Choose one duration $q\in(t,1)$ and split each active endpoint cost
at $z_i=\exp_x(qp_i)$.  The energy inequality
\[
 c(x',y_i)\leq c_q(x',z_i)+c_{1-q}(z_i,y_i)
\]
is an equality at $x'=x$.  Its first term is smooth near $x$, because
$qp_i$ is before cut.  Its gradient is $-p_i$ and its Hessian is
$A_x(qp_i)/q$.  Combining this upper support with the smooth prefix
cost gives
\begin{align}
 A(b-h)&\geq A(b)+hB_o(s_i)
       +\langle p_i-p,\Delta_h\rangle
       +\frac12 Q_i(\Delta_h,\Delta_h)+o(|\Delta_h|^2), \label{sc:individual-pole}\\
 Q_i&=\frac1t A_x(tp)-\frac1q A_x(qp_i).\notag
\end{align}
All expansions here concern fixed scales, a fixed time, and a fixed
finite representation.
Equations~\eqref{sc:pole-movement}--\eqref{sc:individual-pole} bound
each $\langle p_i-p,\Delta_h\rangle$ above by $O(h)$.
Their positive weighted sum is zero, so each is also bounded below by
$-O(h)$.  Hence, for $E=\Span\{p_i-p\}$,
\begin{equation}
 |\operatorname{proj}_E\Delta_h|=O(h).
 \label{sc:linear-movement}
\end{equation}

On $E^\perp$, MTW concavity on the simplex of the $qp_i$ gives
\[
 A_x(qp)\geq\sum_i m_iA_x(qp_i).
\]
Radial monotonicity of divided action, equivalently
Lemma~\ref{geo:slack}, then gives
\begin{equation}
 \sum_i m_iQ_i
 =\frac1t A_x(tp)-\sum_i\frac{m_i}qA_x(qp_i)\geq0
 \quad\text{on }E^\perp.
 \label{sc:transverse-pole}
\end{equation}
Average \eqref{sc:individual-pole}.  The linear terms cancel.
By \eqref{sc:pole-movement}, \eqref{sc:linear-movement}, and
\eqref{sc:transverse-pole}, its quadratic term is bounded below by
$-o(h)$: the transverse part is nonnegative, mixed terms are
$O(h^{3/2})$, and the remaining part is $O(h^2)$.
The Taylor remainders are $o(h)$ as well.  Thus
\[
 A(b-h)-A(b)\geq h\sum_i m_iB_o(s_i)+o(h).
\]
The unique center minimizer and the elementary envelope argument
(compare the old and new minimizers and use their convergence) give
\[
 C(b-h)-C(b)=-hB_c(s)+o(h).
\]
Since $b$ is an interior maximum of $A+C-\pi$, division of the
resulting one-sided comparison by $h$ and passage to the limit give
\eqref{sc:parameter-inequality}.  The argument applies to the
arbitrary representation fixed at the start, proving the asserted
universal quantifier.
\end{proof}

\begin{lemma}[Levels at a regularized maximum]\label{sc:levels}
At a maximum in Lemma~\ref{sc:maximum}, the center and all active
endpoints belong to $S_{(H_o+2)b}(x;u,v)$.
For every active log $p_i$,
\begin{equation}
 v(y_i)-v(y)=\frac{|p|^2-|p_i|^2}{2}
          +b\{B_c(s)-B_o(s_i)\}-E_b(y).
 \label{sc:level-identity}
\end{equation}
In every positive barycentric representation the total weight on
levels $s_i\geq1-\eta$ is at least $c_1$.
Moreover
\begin{equation}
 s<3/4,\qquad
 \max\{s,s_i\}-\min\{s,s_i\}\leq1+\delta_j,\qquad
 \min\{s,s_i\}\geq-2\eta
 \label{sc:level-bounds}
\end{equation}
for small scales.  These conclusions hold at each exact maximum,
before taking a prefix-time limit.
\end{lemma}

\begin{proof}
An active endpoint has original gap
$u-g_b-bB_o(s_i)\in[0,bH_o]$.
By \eqref{sc:aligned}, the center gap is
$u-g_b+E_b(y)-bB_c(s)$, which is at most $(H_o+2)b$ by
\eqref{sc:excess-upper}.  This proves the common section assertion
and its oscillation bound.  Subtracting the active contact equality
from \eqref{sc:aligned} gives \eqref{sc:level-identity}.
Since $E_b(y)>0$ and $E_b(y)\leq bB_c(s)$, one has $B_c(s)>0$,
and therefore $s<3/4$.

Let $m_R$ be the total right weight.  Nonnegativity of $B_o$ and its
lower bound by one outside the right region give
\[
 1-m_R\leq\sum_i m_iB_o(s_i)\leq B_c(s)-c_1\leq1-c_1.
\]
Thus $m_R\geq c_1$.  In particular a right endpoint is present.
The common range bound then puts every level above
$(1-\eta)-(1+\delta_j)=-\eta-\delta_j\geq-2\eta$.
\end{proof}

\subsection{A quantitative outward endpoint}

The next Lemma selects an endpoint once the active logs lie on a line
and the center log has a nonzero component along that line.
The center Hessian estimate in Section~\ref{sec:jets} will establish
these two properties; the selection itself is finite-dimensional.

\begin{lemma}[Weight, cap, and norm increment]\label{sc:outward}
Suppose $p=\sum_{i=1}^N m_ip_i$ in an $n$-dimensional Euclidean
space, with $N\leq n+1$, $m_i>0$, and $\sum_i m_i=1$.
Suppose $E=\Span\{p_i-p\}=\R e$, where $|e|=1$, and
$q=p\cdot e>0$.
Assume the template and level inequalities
\eqref{sc:parameter-inequality}, \eqref{sc:level-bounds}, and
\[
 \left|\ |p_i|^2-|p|^2+2D(s_i-s)\ \right|\leq\zeta D,\qquad
 0\leq\delta_j\leq\zeta\leq1/16.
\]
There is an outward index $i$ with $p_i=p+d_i e$, $d_i>0$, such that
\begin{gather}
 m_i\geq m_*:=\frac{c_1}{12(n+1)},\qquad
 m_id_i\geq\frac{c_1D}{8(n+1)q},\qquad
 d_i\leq\frac{3D}{2q},                         \label{sc:outward-weight}\\
 0<|p_i|^2-|p|^2\leq3D,                       \label{sc:outward-norm}\\
 -2\eta\leq s_i<1-\eta,\qquad B_o(s_i)\leq1/m_*.
                                                        \label{sc:outward-cap}
\end{gather}
In particular the cap $M_0=2/m_*$ can be fixed before $K$.
All these constants are independent of $K$.
\end{lemma}

\begin{proof}
Write $p_i=p+d_i e$.  The level inequality gives total right mass
at least $c_1$, as in Lemma~\ref{sc:levels}.
A right point has $s_i-s>1/4-\eta$, hence
\[
 |p_i|^2-|p|^2=2qd_i+d_i^2\leq-D/4.
\]
It follows that $d_i<0$ and $|d_i|\geq D/(8q)$.
Every outward point satisfies, by the common level range,
\[
 0<2qd_i+d_i^2\leq2D(1+\delta_j)+\zeta D\leq3D,
\]
and so $d_i\leq3D/(2q)$.
The barycenter identity $\sum_i m_id_i=0$ therefore implies
\[
 \sum_{d_i>0}m_id_i\geq\frac{c_1D}{8q}.
\]
Choose an outward index contributing at least $1/N$ of this moment.
Its upper displacement bound gives $m_i\geq c_1/(12N)\geq m_*$,
proving \eqref{sc:outward-weight}.

The parameter inequality and nonnegativity of the outer template give
$m_iB_o(s_i)\leq1$, proving its cap bound.
The lower level bound is already known.  The positive norm increment
implies $s_i\leq s+\zeta/2<3/4+1/32<1-\eta$.
Thus the selected endpoint lies in the region where
\eqref{sc:curvature} applies when $M_0=2/m_*$.
\end{proof}

For a convergent sequence of active representations, an endpoint
selected by Lemma~\ref{sc:outward} in the limit carries a definite
positive moment as well as a definite weight.  Its strict outward
displacement and its strict separation from the right region persist
for the corresponding approximating vertices.  The remaining analytic
task is to produce the
one-dimensional difference space and a bound $q^2\geq\kappa D$ with
$\kappa>0$ independent of $K$. Together with
\eqref{sc:outward-norm}, these are the hypotheses of
Corollary~\ref{spec:outward}, which will compare the center and selected
outer Jacobians.
